%-*-tex-*-
% Copyright Michael J. Ferguson, INRS-Telecommunications
% All rights reserved. 

%======== macros for handling multiple part documents =======
% documents may be broken down to the chapter or section level 


% each part of the document will start with the command 
% \subdocument
% it simply allows for the insertion of special macros or other stuff
% that normally would be brought in only once. In addition, the page numbering
% can be modified at this point to make sense at this point. 

% parameters should be entered to tell whether the sub document is just a
% section, chapter, or a section in a chapter. 
% anything can be set inside the command
% For a chapter (or series) \following\chnum = <chapter number>
%                           \the\chaptername = {<Chapter name>}
%                           \nochaptersfalse
%                           \pagenumber [=] <pagenumber>

%For a section in a document without chapters
%                           \following\shnum = <section number>
%                           \pagenumber [=] <pagenumber>

% for a section in a document with chapters put both together
% it is assumed that the subdocument will start with a section or
% chead. The toplevel will set \subdocumentfalse immediately

% this is here because it is  most useful in sub document context
% multiple documents basically, \following\chnum =2  will set \chnum =1

\def\following#1{\def\a@df{\global\advance #1 by -1\relax} 
                 \afterassignment \a@df \global#1 }

\newif\ifs@ubdocument \s@ubdocumenttrue
\def\subdocumentmacros#1{\ifs@ubdocument #1 \fi\s@ubdocumentfalse}
\outer\def\inputsubdocument #1{\begingroup
                              \s@ubdocumentfalse
                              \let\done=\relax 
                              \input #1 \relax
                              \endgroup}
